\documentclass{article}
\usepackage{amstext}
\usepackage[forcolorpaper]{web}
\usepackage[usesumrytbls,allowrandomize]{exerquiz}[2021/05/29]
\usepackage{ran_toks}[2021/05/29]
\usepackage{multicol}

\DeclareQuiz{Qz1}

\useRandomSeed{1294278923}
%\useLastAsSeed

\hypersetup{pdfpagemode=UseNone} % don't need to see bookmarks
\hypersetup{pdfpagelayout=OneColumn}
\reversemarginpar
\showCreditMarkup

%\previewOn\pmpvOn

\useBeginQuizButton[\CA{Begin}]
\useEndQuizButton[\CA{End}]
\PTsHook{($\eqPTs^{\text{pts}}$)}
\useMCCircles

\newcommand{\presetMatch}{\Q{1}\AddAAKeystroke{event.change=event.change.toUpperCase();}\rectW{\widthof{AA}}}

\begin{document}

\bRTVToks{\currQuiz}

\begin{quiz*}{\currQuiz}
Solve each of these problems, passing is 100\%.
\begin{questions}
    \begin{rtVW}
        \item\PTs{2} $9+8=\RespBoxMath{17}*{1}{.0001}{[0,1]}\cgBdry[1bp]
      \CorrAnsButton{17}$
        \begin{solution}
            Blimey! Everyone knows that $ 9 + 8 = 17$, why don't you ? BTW, $17$ is my formerly favorite number.
\end{solution}
\end{rtVW}

    \begin{rtVW}
        \item\PTs{3} Which of the following are numbers?
        \begin{manswers}*{6}
            \bChoices[random=true]
            \Ans[-1]{0}D\eAns
            \Ans[1]{1}17\eAns
            \Ans[-1]{0}P\eAns
            \Ans[1]{1}88\eAns
            \Ans[-1]{0}S\eAns
            \Ans[1]{1}105\eAns
            \eChoices
        \end{manswers}
        \begin{solution}
         Which are  numbers? 17, 88, and 105, silly
\end{solution}
\end{rtVW}

    \begin{rtVW}
  \item\PTs{3} $ \cos(\pi) = \RespBoxMath{-1}*{1}{.0001}{[0,1]}\cgBdry[1bp]
      \CorrAnsButton{-1} $
\begin{solution}
$ \cos(\pi) = -1 $
\end{solution}
\end{rtVW}

    \begin{rtVW}
\item\PTs{4} $\displaystyle\frac{d}{dx}{\sin(x)}=\RespBoxMath{cos(x)}*{4}{.0001}{[0,1]}\cgBdry[1bp]
      \CorrAnsButton{cos(x)} $
\begin{solution}
 $\displaystyle\frac{d}{dx}{\sin(x)}= \cos(x) $
\end{solution}
\end{rtVW}

    \begin{rtVW}
\hideCreditMarkup
        \item\PTs{3} Match the last name of the U.S. President with his first
        name in the right most two columns. Each problem is worth 1 point.

\bRTVToks{RandomQzQuesB}
\begin{rtVWi}
   \PTs*{1}\RespBoxTxt[\presets{\presetMatch}\rectW{5mm}\MaxLen{1}]{0}{0}*{1}{\txtRef{willy}} Brandt
   \begin{solution}
       \underbar{Willy} Brandt
\end{solution}
\end{rtVWi}
\begin{rtVWi}
   \PTs*{1}\RespBoxTxt[\presets{\presetMatch}\rectW{5mm}\MaxLen{1}]{0}{0}*{1}{\txtRef{roman}} Herzog
   \begin{solution}
       \underbar{Roman} Herzog
\end{solution}
\end{rtVWi}
\begin{rtVWi}
   \PTs*{1}\RespBoxTxt[\presets{\presetMatch}\rectW{5mm}\MaxLen{1}]{0}{0}*{1}{\txtRef{konrad}} Adenauer
   \begin{solution}
       \underbar{Konrad} Adenauer
\end{solution}
\end{rtVWi}
\eRTVToks
\bRTVToks{RandomQzAltsB}
\begin{rtVWi}
   Gustav
\end{rtVWi}
\begin{rtVWi}
   \label{konrad} Konrad
\end{rtVWi}
\begin{rtVWi}
   \label{willy} Willy
\end{rtVWi}
\begin{rtVWi}
   \label{roman} Roman
\end{rtVWi}
\begin{rtVWi}
   Ludwig
\end{rtVWi}
\begin{rtVWi}
   Hermann
\end{rtVWi}
\begin{rtVWi}
   Karl
\end{rtVWi}
\begin{rtVWi}
   J\"{u}rgen
\end{rtVWi}
\eRTVToks
    \begin{multicols}{4}
       \begin{questions}[labelwidthTo={\textbf{A}}]
           \hideCreditMarkup
           \renewcommand{\theeqquestionnoii}{\Alph{eqquestionnoii}}
           \displayListRandomly[\item]{RandomQzAltsB}
           \showCreditMarkup
       \end{questions}
    \end{multicols}
    \begin{questions}
       \showCreditMarkup
       \displayListRandomly[\item]{RandomQzQuesB}
    \end{questions}
\end{rtVW}
  \eRTVToks
  \displayListRandomly{\thisQuiz}
\end{questions}
%\writeProListAux
\end{quiz*}\cgBdry[1em]\ScoreField{\currQuiz}\olBdry\CorrButton{\currQuiz}\vcgBdry[6pt]

Answers: \AnswerField\currQuiz

\begin{center}
\displaySumryTbl[showmarkup]{\currQuiz}
\end{center}
\end{document}
